\chapter{Оршил}

Энэхүү бакалаврын судалгааны ажлын зорилго нь хөдөлгөөнт төхөөрөмжөөс авсан дэлгүүрийн лангууны зурагт гүн сургалтын объект илрүүлэлтийн аргыг ашиглан бараа бүтээгдэхүүнийг автоматаар таньж, илрүүлэлтийн үр дүнд тулгуурласан аудитын шийдвэр гаргалтыг автоматжуулах төгсгөл хоорондын систем боловсруулахад оршино. Судалгааны төв асуудал нь зөвхөн ``зургаас бараа олох'' бус, харин дүрсийн боловсруулалт, серверийн боловсруулалт, өгөгдөл хадгалалт, аудитын дүрэм, хэрэглэгчийн урсгалыг нэгэн архитектурт нэгтгэхэд оршиж байна.

\section{Зорилтууд}

Дээрх ерөнхий зорилгыг хэрэгжүүлэхийн тулд дараах тодорхой зорилтуудыг дэвшүүлэв:

\begin{enumerate}
\item \textbf{Онолын судалгаа хийх:} Компьютерын харааны объект илрүүлэлтийн алгоритмуудын хөгжлийн чиг хандлага, архитектурыг судлах, жижиглэн худалдааны салбарт ашиглагдаж буй автоматжуулалтын системүүдийн давуу болон сул талыг харьцуулан шинжлэх.
\item \textbf{Өгөгдөл бүрдүүлж, загвар сургах:} Бодит орчинд лангууны барааны зургийн өгөгдлийн санг цуглуулж (300+ зураг), объект бүрийг тэмдэглэх замаар датасет бэлтгэх. YOLOv8 архитектурыг ашиглан 8 ангиллын барааг таних гүн сургалтын загвар сургах.
\item \textbf{Сервер хөгжүүлэх:} FastAPI фреймворк ашиглан зураг хүлээн авч, сургасан загвараар объектуудыг илрүүлэх, MongoDB өгөгдлийн санд үр дүнг хадгалах өндөр хурдны асинхрон REST API хөгжүүлэх.
\item \textbf{Автомат аудитын алгоритм боловсруулах:} Зургаас таньсан үр дүнг хүлээгдэж буй барааны мэдээлэлтэй харьцуулан зөрүүг тооцоолж, аудитын статус (PASS, WARNING, FAIL) гаргах шийдвэр гаргалтын логик боловсруулах.
\item \textbf{Гар утасны аппликейшн хөгжүүлэх:} Flutter технологи ашиглан iOS болон Android дээр ажиллах, хэрэглэгчдэд ээлтэй мобайл аппликейшн бүтээх.
\end{enumerate}

\section{Судалгааны хэрэгцээ шаардлага}

Дэлхий даяар жижиглэн худалдааны салбар нь эдийн засгийн томоохон хэсгийг бүрдүүлдэг. Statista-ийн 2023 оны тайлангаар дэлхийн жижиглэн худалдааны зах зээлийн хэмжээ 28.9 их наяд ам.доллар хүрсэн бөгөөд 2027 он гэхэд 32.8 их наяд ам.долларт хүрэхээр таамаглагдаж байна \cite{statista_retail}. Ийм хэмжээний зах зээлд лангууны дүүргэлт, барааны олдоц, тооллогын найдвартай байдал нь зөвхөн үйл ажиллагааны бус, орлогын түвшний асуудал болдог.

McKinsey \& Company-ийн 2023 оны судалгаагаар жижиглэн худалдааны компаниудын 67\% нь бараа материалын менежментийг сайжруулахын тулд хиймэл оюун ухааны технологи нэвтрүүлэхэд бэлэн байгаагаа илэрхийлсэн \cite{mckinsey_retail}. Гэвч өндөр өртөг, дэд бүтцийн шаардлага, нарийн мэргэжлийн нөөцийн хомсдол нь ийм шийдлийг өргөн хэрэглэхэд саад болж байна.

Монгол Улсын хувьд асуудлын орон нутгийн ач холбогдол мөн өндөр байна. Үндэсний Статистикийн Хорооны мэдээллээр дотоод худалдааны салбарын жижиглэн худалдааны нийт борлуулалт 2020 онд 4{,}821.3 тэрбум төгрөг байсан бол 2024 онд 7{,}520.0 тэрбум төгрөгт хүрч, таван жилийн хугацаанд тогтвортой өссөн байна \cite{nso_mongolia}. Энэ өсөлт нь супермаркет, convenience store, сүлжээ дэлгүүрийн орчинд бодит цагийн нөөцийн хяналт, лангууны мониторинг, аудитын автоматжуулалтын хэрэгцээ нэмэгдэж байгааг илтгэнэ.

\begin{longtable}{|p{2.2cm}|p{4cm}|p{4cm}|}
\caption{Монгол Улсын жижиглэн худалдааны борлуулалтын өсөлт} \label{tab:intro_mn_retail} \\
\hline
\textbf{Он} & \textbf{Нийт борлуулалт} & \textbf{Жилийн өөрчлөлт} \\
\hline
\endfirsthead
\hline
\textbf{Он} & \textbf{Нийт борлуулалт} & \textbf{Жилийн өөрчлөлт} \\
\hline
\endhead
\hline
\endfoot
2020 & 4{,}821.3 тэрбум төгрөг & --- \\
\hline
2021 & 5{,}234.7 тэрбум төгрөг & 8.6\% өсөлт \\
\hline
2022 & 6{,}012.4 тэрбум төгрөг & 14.9\% өсөлт \\
\hline
2023 & 6{,}845.1 тэрбум төгрөг & 13.8\% өсөлт \\
\hline
2024 & 7{,}520.0 тэрбум төгрөг & 9.9\% өсөлт \\
\hline
\end{longtable}

Энэхүү өрсөлдөөнт орчинд бараа материалын менежмент, ялангуяа лангуу дээрх өрөлт, тооллого, татан авалтын хяналт онцгой ач холбогдолтой. IHL Group-ийн ``Retailers and the Ghost Economy'' тайлангаар бараа дууссанаас үүдэлтэй борлуулалтын алдагдал жил бүр 1.75 их наяд ам.долларт хүрдэг бөгөөд энэ нь жижиглэн худалдаачдын нийт орлогын ойролцоогоор 4\%-ийг эзэлдэг \cite{retail_stats}. Үүний нэг хэсэг нь лангуунд бараа бодитоор байсан ч системд буруу бүртгэгдсэн эсвэл зөв байрлалд байгаагүйгээс үүсэх ``phantom inventory'' асуудалтай холбоотой.

\begin{table}[H]
\centering
\caption{Жижиглэн худалдааны салбарын бараа материалын алдагдлын статистик}
\label{tab:retail_loss_stats}
\begin{longtable}{|p{5cm}|c|c|}
\hline
\textbf{Алдагдлын төрөл} & \textbf{Жилийн алдагдал (ам.доллар)} & \textbf{Эзлэх хувь} \\
\hline
Бараа дууссанаас үүдсэн & 1.75 их наяд & 53\% \\
\hline
Буруу байршуулалт & 634 тэрбум & 19\% \\
\hline
Хугацаа дууссан бараа & 408 тэрбум & 12\% \\
\hline
Тооллогын алдаа & 531 тэрбум & 16\% \\
\hline
\textbf{Нийт} & \textbf{3.32 их наяд} & \textbf{100\%} \\
\hline
\end{longtable}
\end{table}

Уламжлалт жижиглэн худалдааны орчинд дэлгүүрийн ажилтан эсвэл борлуулалтын төлөөлөгч лангууг биечлэн шалгаж, дууссан болон дутуу өрөгдсөн барааг гараар бүртгэдэг. Энэ ажиллагаа нь хугацаа их шаарддагаас гадна анхаарал сарних, нэгж алгасах, ижил төрлийн савлагаа андуурах зэрэг алдаа дагуулдаг. Zebra Technologies-ийн 2023 оны судалгаагаар жижиглэн худалдааны ажилчдын 72\% нь лангууны тооллогод хэт их цаг зарцуулдаг гэж хариулсан бөгөөд 35\% нь тооллогын үр дүнг найдваргүй гэж үзсэн \cite{zebra_study}.

Сүүлийн жилүүдэд компьютерын хараа, гүн сургалтын хөгжил ийм асуудлыг зурагт суурилсан аргаар шийдвэрлэх боломжийг бодит түвшинд ойртуулсан. Grand View Research-ийн 2024 оны тайлангаар жижиглэн худалдаанд хэрэглэгддэг компьютерын харааны зах зээл 2023 онд 2.8 тэрбум ам.доллар байсан бөгөөд 2030 он гэхэд 13.5 тэрбум ам.долларт хүрч жилийн дундаж өсөлт 25.3\% байхаар таамаглагдаж байна \cite{grand_view_cv}. Иймээс энэхүү судалгаа нь зөвхөн алгоритмын туршилт бус, мобайл зураг цуглуулалт, серверийн inference, аудитын шийдвэр гаргалтыг уялдуулсан хэрэглээний судалгаа болно.

\section{Тулгамдсан асуудлууд}

Бараа материалын менежмент нь бодит худалдааны орчинд олон давхар хязгаарлалттай тулгардаг. Лангуун дээр нягт байрласан ижил төрлийн савлагаа, өрөлтийн байнгын өөрчлөлт, гэрэлтүүлгийн хэлбэлзэл, камерын өнцгийн ялгаа зэрэг хүчин зүйлс нь автомат илрүүлэлтийн найдвартай байдалд шууд нөлөөлдөг \cite{real_time_detection,sku110k,retail_cv_survey}. Ийм орчинд гараар хийдэг шалгалт удаан төдийгүй тогтвортой нарийвчлалтай байж чаддаггүй.

Монголын нөхцөлд асуудлын эдийн засгийн хэмжээг ойролцоогоор тооцож болно. Хэрэв олон улсад өргөн эшлэгддэг 4.1\%-ийн борлуулалтын алдагдлын логикийг \cite{mckendrick_ai} Монгол Улсын 2024 оны 7{,}520.0 тэрбум төгрөгийн жижиглэн худалдааны эргэлтэд хэрэглэвэл:
\begin{equation}
Loss_{MN} = 7{,}520.0 \times 0.041 = 308.3 \text{ тэрбум төгрөг}
\end{equation}
гэж тооцоологдоно. Энэ тооцоо нь Монголын бүх жижиглэн худалдааны байгууллагын бодит хохирлыг шууд хэмжсэн үзүүлэлт биш боловч нөөцийн зөрүү, барааны дутагдал, тооллогын алдаа зэрэг асуудал макро түвшинд ямар хэмжээний эдийн засгийн эрсдэл дагуулж болохыг харуулах аналитик үндэслэл юм. Иймээс автомат илрүүлэлт ба аудитын систем нь зөвхөн техникийн шинэчлэл бус, зардлын хяналтын хэрэгсэл мөн болно.

Ийм нөхцөлөөс дараах үндсэн асуудлууд тодорхойлогдоно:

\begin{enumerate}
\item \textbf{Хүний нөөцийн зардал ба цаг хугацааны алдагдал:} Олон нэр төрлийн барааг тогтмол шалгах нь мэдэгдэхүйц хөдөлмөр, цаг хугацаа шаарддаг тул аудитын давтамж, хамрах хүрээг бууруулах эрсдэлтэй.
\item \textbf{Хүний хүчин зүйлийн алдаа:} Давтагдсан тооллогын үед анхаарал сарних, ижил савлагаа андуурах, нэгж алгасах зэрэг алдаа өсдөг бөгөөд энэ нь тайлангийн найдвартай байдалд шууд нөлөөлнө.
\item \textbf{Бодит хугацааны мэдээллийн саатал:} Шалгалтын үр дүнг төв системд оруулах хүртэл хоцролт үүсэх тул нөхөн дүүргэлт, нийлүүлэлтийн шийдвэр саатах эрсдэлтэй.
\item \textbf{Дэд бүтцийн өндөр өртөг:} Зарим автоматжсан шийдлүүд нь тусгай камер, мэдрэгч, иж бүрэн суурилуулалт шаарддаг бөгөөд энэ нь жижиг, дунд байгууллагад нэвтрүүлэх өртгийг өсгөдөг \cite{amazon_go}.
\item \textbf{Орон нутгийн өгөгдөлд дасан зохицох шаардлага:} Олон улсын датасет ба бэлэн загварууд нь Монголын шошго, савлагаа, өрөлтийн хэв маягтай бүрэн нийцэхгүй тул орон нутгийн fine-tuning шаарддаг \cite{transfer_learning,retail_cv_survey}.
\end{enumerate}

Эдгээр асуудал нь зөвхөн detection accuracy-ийн биш, харин бизнесийн үр ашиг, шийдвэр гаргалтын хурд, бодит орчинд нэвтрүүлэх боломжийн нийлмэл асуудал юм. Иймээс зурагт суурилсан, хямд өртөгтэй, орон нутгийн өгөгдөлд дасан зохицсон автомат аудитын шийдэл боловсруулах нь судалгааны болон практикийн аль алинд ач холбогдолтой.

\section{Судалгааны гол хувь нэмэр}

Сүүлийн үеийн тойм болон салбарын судалгаанууд нь бараа материалын менежмент, мобайл зураг цуглуулалт, компьютерын хараанд тулгуурласан бүтээгдэхүүн танилт бодит хэрэглээнд шилжих боломжтой түвшинд хүрснийг харуулж байна \cite{inventory_management_ai,mobile_retail,retail_cv_survey}. Гэсэн хэдий ч жижиг, дунд худалдааны байгууллагад зориулсан, хямд өртөгтэй бөгөөд төгсгөл хоорондын нэгтгэлтэй шийдэл харьцангуй цөөн хэвээр байна. Энэхүү судалгааны гол хувь нэмэр нь энэ орон зайг нөхөхөд чиглэнэ.

Энэхүү судалгааны ажлын гол хувь нэмрүүд дараах байдалтай:

\begin{enumerate}
\item \textbf{Төгсгөл хоорондын автомат аудитын архитектур:} Зураг авах, бүтээгдэхүүн илрүүлэх, тоо нэгтгэх, нормтой харьцуулах, аудитын статус гаргах дамжлагыг нэг урсгалд холбосон.
\item \textbf{Орон нутгийн судалгааны датасет:} Улаанбаатар хотын жижиглэн худалдааны орчноос цуглуулсан 324 зураг, 2{,}847 annotation бүхий өгөгдөл нь Монголын савлагаа, шошго, өрөлтийн хэв маягийг тусгасан domain-specific fine-tuning-ийн суурь болсон.
\item \textbf{Орон нутгийн нөхцөлд дасан зохицсон загварчлал:} YOLOv8 суурьтай detector-ийг локал өгөгдөлд тохируулснаар олон улсын ерөнхий загварын domain shift-ийг бууруулах боломжийг харуулсан \cite{transfer_learning,retail_cv_survey}.
\item \textbf{Хэрэглээний хувьд хөнгөн шийдэл:} Ухаалаг утас ба хөнгөн сервер дээр ажиллах зохиомж санал болгосноор тусгай мэдрэгч, өндөр өртөгтэй дэд бүтэц шаардахгүй байх боломжийг харуулсан.
\item \textbf{Нээлттэй, өргөтгөх боломжтой технологийн суурь:} YOLOv8, FastAPI, MongoDB, Flutter зэрэг технологийг уялдуулснаар системийг турших, өргөтгөх, дахин ашиглах инженерийн суурь бүрдүүлсэн.
\end{enumerate}

Ингэснээр оршил бүлэгт тодорхойлсон хэрэгцээ, асуудал, зорилт, хувь нэмэр нь дараагийн бүлгүүдийн судалгааны тойм, системийн зохиомж, хэрэгжүүлэлт, үнэлгээтэй шууд уялдах бөгөөд төгсгөлд зорилт тус бүрийн биелэлтийг тоон үр дүнгээр хаах бүтэц бүрдэнэ.

\section{Судалгааны объект ба хамрах хүрээ}

Энэхүү судалгааны объект нь дэлгүүрийн лангууны зураг дээр суурилсан бүтээгдэхүүн илрүүлэлт болон түүнээс үүсэх аудитын шийдвэр гаргалтын үйл явц юм. Судалгааны төв цөм нь дүрсээс мэдээлэл гарган авах, түүнийг бизнесийн дүрэмтэй холбох pipeline-д оршино. Хамрах хүрээний хувьд дараах хязгаарлалтыг тогтоов:

\begin{itemize}
    \item Туршилтад өргөн хэрэглээний барааны ангиллуудыг хамарсан зураг ашиглаж, лангууны орчин дахь object detection-ийн гүйцэтгэлийг үнэлсэн.
    \item Систем нь интернетэд холбогдсон гар утаснаас зураг илгээж, сервер талд inference хийх загварт тулгуурлана.
    \item Судалгаа нь object detection, тооллого, аудитын статус тодорхойлолтод төвлөрөх бөгөөд гар утсан дээрх шууд edge inference-ийг үндсэн зорилгоос тусад нь авч үзэв.
    \item Барааны үнэ, штрихкод, кассын түвшний бүртгэл зэрэг нэмэлт модулиудыг цаашдын өргөтгөл гэж үзсэн.
\end{itemize}

% \section{Судалгааны ач холбогдол}

% \subsection{Шинжлэх ухааны ач холбогдол}

% Компьютерын харааны object detection нь нягт байрласан, хэсэгчлэн халхлагдсан орчинд оновчтой шийдэл сонгохыг шаарддаг төвөгтэй асуудал хэвээр байна \cite{real_time_detection,sku110k}. Энэ судалгаа нь ийм нөхцөлд YOLOv8 суурьтай detector-ийг бодит аудитын урсгалтай уялдуулан туршсанаараа онолын болон аргазүйн ач холбогдолтой.

% \begin{itemize}
%     \item Орчин үеийн detector-ийг барааны лангуу шиг нягт дүрслэлтэй орчинд ашиглах боломжийг практик туршилтаар харуулсан.
%     \item Илрүүлэлтийн үр дүнг аудитын шийдвэртэй холбосон төгсгөл хоорондын архитектурын жишиг загвар санал болгосон.
%     \item Компьютерын хараа, мобайл зураг цуглуулалт, сервер талын боловсруулалтыг нэгтгэсэн системийн зохиомжийн судалгаанд хувь нэмэр оруулсан.
% \end{itemize}

% \subsection{Практикийн ач холбогдол}

% Практикийн хувьд энэхүү шийдэл нь дэлгүүрийн ажилтны гар ажиллагааг багасгаж, лангуун дээрх барааны байршил, тоо ширхэгийн талаарх мэдээллийг илүү хурдан боловсруулах боломж олгоно. Ухаалаг утас болон хөнгөн серверийн дэд бүтэц дээр хэрэгжих боломжтой тул өргөн хэрэглээнд нэвтрэх босго харьцангуй бага \cite{mobile_retail,inventory_management_ai}.

% \begin{itemize}
%     \item Бараа тооллого болон аудитын бэлтгэлд зарцуулах гар ажиллагааг бууруулж, хүний нөөцийн ачааллыг хөнгөлнө.
%     \item Удирдлагад бараа дутагдал, зөрүү, өрөлтийн алдааг илүү богино хугацаанд илрүүлэх боломж олгоно.
%     \item Нэмэлт тусгай тоног төхөөрөмж шаардахгүй тул жижиг, дунд худалдааны байгууллагад ашиглах боломжтой хямд өртөгтэй сонголт болно.
%     \item Цаашид өөр барааны ангилал, шинэ дэлгүүрийн формат, өөр үйлчилгээний загварт өргөтгөх суурь архитектур бүрдүүлнэ.
% \end{itemize}
